\section*{Chapitre \ref{ch:g12:polymers} --- Les polymères}

\begin{solution}{exo:g12:polymers:1}
Le propène, \ce{CH2=CH-CH3}.
\end{solution}

\begin{solution}{exo:g12:polymers:2}
\ce{-CF2-CF2-}~; un \omterm{def:g12:polymers:addition-polymer}{polymère d'addition}~: le \omterm{def:g12:polymers:polymer}{monomère} possède une
liaison \ce{C=C} et il ne se forme aucun autre produit.
\end{solution}

\begin{solution}{exo:g12:polymers:3}
\omterm{def:g12:curly-arrows:categories}{Addition}~: le polyéthylène, le PVC, le polystyrène. Condensation~: le
PET (un polyester), le nylon (un polyamide).
\end{solution}

\begin{solution}{exo:g12:polymers:4}
La cellulose, l'amidon, le caoutchouc issu du latex, les protéines. Le
nylon et le PVC sont synthétiques.
\end{solution}

\begin{solution}{exo:g12:polymers:5}
Chaque liaison utilise un groupe de chaque \omterm{def:g12:polymers:polymer}{monomère}~; avec deux groupes,
un \omterm{def:g12:polymers:polymer}{monomère} peut se lier des deux côtés et la chaîne continue de
croître. Un \omterm{def:g12:polymers:polymer}{monomère} à un seul groupe arrêterait la chaîne.
\end{solution}

\begin{solution}{exo:g12:polymers:6}
$M(\ce{C2H3Cl}) = 24.0 + 3.0 + 35.5 = \qty{62.5}{g/mol}$;
$n = 125000/62.5 = 2000$.
\end{solution}

\begin{solution}{exo:g12:polymers:7}
$M(\ce{C8H8}) = \qty{104.0}{g/mol}$; $M = 2500 \times 104.0 =
\qty{2.60e5}{g/mol}$.
\end{solution}

\begin{solution}{exo:g12:polymers:8}
Le groupe acide d'une \omterm{def:g7:atoms-and-molecules:molecule}{molécule} forme un \omterm{def:g11:functional-groups:carboxylic-acid}{ester} avec le groupe \omterm{def:g11:functional-groups:alcohol}{alcool} de
la suivante, qui possède encore un groupe acide libre~: la chaîne
croît. \omterm{def:g12:polymers:repeat-unit}{Motif}
\ce{-O-CH(CH3)-CO-}~; de l'eau est libérée à chaque liaison.
\end{solution}

\begin{solution}{exo:g12:polymers:9}
Les films pour sacs sont faits de chaînes ramifiées, qui ne peuvent pas
se ranger~: le matériau est amorphe, souple et transparent. Les
bouteilles de lait sont faites de chaînes presque linéaires, qui se
rangent en régions cristallines~: plus denses, plus rigides, et les
régions cristallines diffusent la lumière.
\end{solution}

\begin{solution}{exo:g12:polymers:10}
Leurs chaînes sont reliées par de nombreux pontages covalents. Le
chauffage ne peut pas faire glisser les chaînes~; il rompt des liaisons
et détruit le matériau au lieu de le faire fondre.
\end{solution}

\begin{solution}{exo:g12:polymers:11}
$M(\ce{C12H22N2O2}) = 144.0 + 22.0 + 28.0 + 32.0 = \qty{226.0}{g/mol}$;
$n = 22600/226.0 = 100$.
\end{solution}

\begin{solution}{exo:g12:polymers:12}
\[
  \ce{H2N(CH2)6NH2 + HOOC(CH2)4COOH -> H2N(CH2)6NHCO(CH2)4COOH + H2O} .
\]
Le produit porte encore un groupe \omterm{def:g11:functional-groups:amine}{amine} libre à une extrémité et un
groupe acide libre à l'autre~: chacun peut réagir avec un autre
\omterm{def:g12:polymers:polymer}{monomère}.
\end{solution}

\begin{solution}{exo:g12:polymers:13}
Deux \omterm{def:g7:atoms-and-molecules:molecule}{molécules} d'eau par \omterm{def:g12:polymers:repeat-unit}{motif}. $n = 1000/226.0 =
\qty{4.42}{mol}$ de \omterm{def:g12:polymers:repeat-unit}{motifs}, donc \qty{8.85}{mol} d'eau~:
$8.85 \times 18.0 = \qty{159}{g}$.
\end{solution}

\begin{solution}{exo:g12:polymers:14}
Le caoutchouc naturel chaud est fait de chaînes qui glissent les unes
sur les autres. Les ponts soufrés relient les chaînes~: elles s'étirent
mais reviennent en place, et ne coulent plus. Étant réticulé, un pneu ne
peut pas être fondu et remis en forme.
\end{solution}

\begin{solution}{exo:g12:polymers:15}
$M(\ce{C6H10O5}) = \qty{162.0}{g/mol}$~; $n = \num{1.6e6}/162.0 = 9900$
\omterm{def:g12:polymers:repeat-unit}{motifs}~; longueur $9900 \times 0.5 = \qty{4900}{nm}$, soit environ
\qty{5}{\micro m}.
\end{solution}

\begin{solution}{pb:g12:polymers:1}
\textbf{1.} L'éthane-1,2-diol (deux groupes \omterm{def:g11:functional-groups:alcohol}{alcool}) et l'acide
benzène-1,4-dicarboxylique (deux groupes \omterm{def:g11:functional-groups:carboxylic-acid}{acide carboxylique}).

\textbf{2.} Une liaison \omterm{def:g11:functional-groups:carboxylic-acid}{ester}~; de l'eau est libérée.

\textbf{3.} Un \omterm{def:g12:polymers:addition-polymer}{polymère de condensation}.

\textbf{4.} $M(\ce{C2H6O2}) = \qty{62.0}{g/mol}$; $M(\ce{C8H6O4}) =
\qty{166.0}{g/mol}$.

\textbf{5.} $62.0 + 166.0 - 2 \times 18.0 = \qty{192.0}{g/mol}$~; et
$10 \times 12.0 + 8 \times 1.0 + 4 \times 16.0 = 192.0$.

\textbf{6.} $n = 38400/192.0 = 200$.

\textbf{7.} Deux liaisons \omterm{def:g11:functional-groups:carboxylic-acid}{ester} par \omterm{def:g12:polymers:repeat-unit}{motif}~: environ 400.

\textbf{8.} Ses chaînes sont surtout linéaires, et certaines de leurs
parties s'alignent en régions ordonnées.

\textbf{9.} C'est un \omterm{def:g8:plastics:thermoplastic}{thermoplastique}~: ses chaînes ne sont pas pontées
et glissent les unes sur les autres quand il est chaud.

\textbf{10.} $300/0.80 = \qty{375}{g}$.

\textbf{11.} $375/25 = 15$ bouteilles.

\textbf{12.} Les bouchons, les étiquettes, les salissures et les chutes
sont retirés ou perdus au cours du tri, du lavage et du filage.

\textbf{13.} Prévenir les déchets (principe 1)~: un objet usagé devient
une \omterm{def:g5:raw-materials-and-recycling:raw-material}{matière première}.

\textbf{14.} \ce{C10H8O4 + 2H2O -> C2H6O2 + C8H6O4}.

\textbf{15.} $1000/192.0 = \qty{5.21}{mol}$.

\textbf{16.} Diol~: $5.21 \times 62.0 = \qty{323}{g}$~; acide~:
$5.21 \times 166.0 = \qty{865}{g}$.

\textbf{17.} $2 \times 5.21 \times 18.0 = \qty{188}{g}$ d'eau~;
$1000 + 188 = 323 + 865 = \qty{1188}{g}$~: la masse est conservée.

\textbf{18.} Les \omterm{def:g12:polymers:polymer}{monomères} peuvent être purifiés et donnent du PET neuf
d'aussi bonne qualité, alors que la fusion raccourcit un peu les chaînes
à chaque fois.

\textbf{19.} Une veste polaire nécessite 15 bouteilles de \qty{25}{g}.
\end{solution}
